% Generated draft: requires mathematical and source review.
\documentclass[11pt]{article}
\usepackage[margin=25mm]{geometry}
\usepackage{fontspec}
\setmainfont{Latin Modern Roman}
\usepackage{amsmath,amssymb,mathtools,unicode-math}
\setmathfont{Latin Modern Math}
\usepackage{enumitem,needspace,xurl,hyperref}
\setlength{\parindent}{0pt}
\setlength{\parskip}{5pt}
\setlength{\emergencystretch}{3em}
\newcommand{\N}{\mathbb{N}}
\newcommand{\Z}{\mathbb{Z}}
\newcommand{\Q}{\mathbb{Q}}
\newcommand{\R}{\mathbb{R}}
\newcommand{\C}{\mathbb{C}}
\newcommand{\F}{\mathbb{F}}
\newcommand{\A}{\mathbb{A}}
\newcommand{\PP}{\mathbb{P}}
\newcommand{\E}{\mathbb{E}}
\newcommand{\eps}{\varepsilon}
\newcommand{\dd}{\,\mathrm{d}}
\begin{document}
\section*{UnsolvedMath ID: 701561}
\textbf{Source problem code:} catalogue entry
\subsection{Definitions and mathematical statement}
A group \(G\) is periodic if every element has finite order. For a prime \(p\), a \(p\)-group means a group in which every element has order a power of \(p\); such a group need not be finite. Write \(a^g=g^{-1}ag\).

For an odd prime \(p\), a periodic group \(G\), and \(a\in G\), consider the hypotheses
\[
\begin{aligned}
(H_{\mathrm{fin}})\quad&\text{for all }g,h\in G,\quad \langle a^g,a^h\rangle\text{ is a finite }p\text{-group};\\
(H_p)\quad&\text{for all }g,h\in G,\quad \langle a^g,a^h\rangle\text{ is a }p\text{-group}.
\end{aligned}
\]
The desired conclusion is
\[
\exists N\trianglelefteq G\quad\bigl(a\in N\ \text{and }N\text{ is a }p\text{-group}\bigr).
\]
The supplied question asks whether the Baer--Suzuki implication holds for every odd prime \(p\), every periodic group \(G\), and every \(a\in G\). Retaining the explicit finiteness condition from the supplied finite-group theorem gives the implication from \((H_{\mathrm{fin}})\) to the displayed conclusion. Allowing the two-conjugate subgroups to be infinite gives the implication from \((H_p)\) to that conclusion. The source does not explicitly distinguish these extensions, so the intended hypothesis requires review. No local finiteness or bounded exponent assumption is supplied. The conclusion above allows \(N\) to be infinite; the source does not expressly discuss this distinction in the periodic setting.

The supplied record states that a counterexample is known for \(p=2\), referring to Archive 11.11a). This is retained as a source assertion, without independent verification.
\subsection{Short English statement}
Borovik asks whether the Baer--Suzuki criterion extends to periodic groups for odd primes: must an element whose every two conjugates generate a suitable \(p\)-subgroup belong to a normal \(p\)-subgroup? The source leaves implicit whether the generated subgroups must remain finite in this extension. The registry labels the problem open; this is a supplied status, not an independently verified assessment.
\subsection{Sources}
\begin{itemize}
\item Supplied local UnsolvedMath registry record, permanent ID \(701561\), exact source code \texttt{catalogue entry}, titled \emph{Kourovka Notebook Problem 19.13}. The record attributes the problem to A. V. Borovik, proposed in 2018, and labels it open.
\item \emph{The Kourovka Notebook (Unsolved Problems in Group Theory)}, Problem 19.13, 19th issue (2018); the supplied background also identifies the 21st issue edition (2026), \texttt{arXiv:1401.0300v48}. The supplied collection description names E. I. Khukhro and V. D. Mazurov as editors.
\item Archive 11.11a), cited by the supplied statement for a counterexample when \(p=2\). The archive entry and details of the counterexample were not supplied.
\end{itemize}
\end{document}
